% D.75. Root and FGH independent complete proof/source/certificate audits passed.
\subsection{Explicit free factors in the retained compact obstruction}

\begin{statement}{D.75}{Proved}{Local Kummer divisibility produces a compact global obstruction factor}
Use the actual Alexander action of D.35 and the compact kernel
\(\mathcal K_s\) in D.63. Let
\[
 0\ne a=\sum_p a_p\kappa(p)\in H^1(G_{\mathbb Q},\mathbb Z_5(1)),
 \qquad s=(a,0).
\]
The coefficients \(a_p\) tend to zero; finite support is not required.
Put
\[
 \begin{split}
 r&=\min_p v_5(a_p),\\
 N&=\sup\{n\ge0:a|_{G_{\mathbb Q_5}}\in
                   5^nH^1(G_{\mathbb Q_5},\mathbb Z_5(1))\},
 \qquad v_5(0)=+\infty .
 \end{split}                                               \tag{D.75.1}
\]
Thus \(r<\infty\), \(N\ge r\), and \(N=\infty\) means that the
local Kummer class is zero. There is a continuous surjection
\[
 \mathcal K_s\twoheadrightarrow
 \begin{cases}
       \mathbb Z_5^{\,N-r},&N<\infty,\\
       \displaystyle\prod_{\mathbb N}\mathbb Z_5,&N=\infty .
 \end{cases}                                               \tag{D.75.2}
\]
It splits continuously as a map of compact \(\mathbb Z_5\)-modules,
noncanonically. In particular these are actual direct factors of the
compact term, not merely classes at one finite coefficient level.

Both integers in (D.75.1) are explicit. With the usual convergent
five-adic logarithm set
\[
 \lambda_p=\frac{\log(p^4)}{4\log 6}\quad(p\ne5),\qquad
 b_5(a)=\sum_{p\ne5}a_p\lambda_p.
\]
Then
\[
                    N=\min\{v_5(a_5),v_5(b_5(a))\}.         \tag{D.75.3}
\]
The displayed sum converges, independently of order.
For example \(s=(\kappa(101),0)\) gives a split
\(\mathbb Z_5\) factor in \(\mathcal K_s\). More generally a prime
\(q\) with \(v_5(q-1)=N+1\), \(N\ge1\), gives a split
\(\mathbb Z_5^N\) factor.

There is also a finite-support example with an infinite free factor.
Put
\[
 c=\frac{\log 16}{\log 81}\in\mathbb Z_5^\times,\qquad
 a_*=\kappa(2)-c\kappa(3).
                                                               \tag{D.75.4}
\]
Here \(c\equiv3\pmod5\), \(r=0\), and \(N=\infty\).
Consequently \(\mathcal K_{(a_*,0)}\) has a continuous direct
factor \(\prod_{\mathbb N}\mathbb Z_5\). Both pairs in these
examples are cusp pairs and have actual sections.
The same assertions hold after permuting the three punctures.

This evaluates a free factor of the compact term in Task D.5.
The complementary compact factor and the general first
finite-derived obstruction are not determined here.
In particular the compact term cannot in general be discarded:
the zero calculation of D.64 does not extend to all finite
ramification sets.
\end{statement}

\begin{proof}
\paragraph{An actual one-variable quotient and its measure model.}
Write \(t=1+u\). The ideal \(vA_s\) is stable. In D.35's
normalization, the integral beta series is one when its second
argument is one. Thus on the actual quotient
\(M_a=A_s/vA_s\simeq\mathbb Z_5[[u]]\) the action is
\[
 g f(t)=t^{a(g)}\chi(g)S_{\chi(g)}(t)f(t^{\chi(g)}),
 \qquad S_\alpha(t)=\frac{t^\alpha-1}{t-1}.
                                                               \tag{D.75.5}
\]
This uses the exact Fox multiplier, including its factor
\(S_{\chi}(1)=\chi\), not a leading-weight approximation.
It follows directly from D.35 and D.39.
Multiplication by \(u\) is an isomorphism of compact modules
\[
 M_a(-1)\xrightarrow{\ \sim\ }I,\qquad
 I=\ker\{\mathbb Z_5[[\mathbb Z_5]]\longrightarrow\mathbb Z_5\},
 \qquad f(t)\longmapsto(t-1)f(t).                            \tag{D.75.6}
\]
On the right the action is the actual affine pushforward
\[
                  \nu\longmapsto t^{a(g)}\phi_{\chi(g)}(\nu),
 \qquad \phi_\alpha(t^z)=t^{\alpha z}.
                                                               \tag{D.75.7}
\]
Indeed \((t-1)S_\chi(t)=t^\chi-1\).
An integral change of Kummer cocycle by a coboundary corresponds
to conjugating this action by the translation \(t^d\), \(d\in\mathbb Z_5\).
Translation preserves \(I\).

Let \(D_a=M_a^D=\operatorname{Hom}_{\rm cts}(M_a,\mathbb Q_5/\mathbb Z_5)(1)\).
The dual of (D.75.6) identifies it with
\[
 D_a\simeq
 \frac{\operatorname{Cont}(\mathbb Z_5,\mathbb Q_5/\mathbb Z_5)}
                         {\mathbb Q_5/\mathbb Z_5}.
                                                               \tag{D.75.8}
\]
The denominator consists of constant functions. To verify this
identification integrally, a continuous linear functional on the
completed group ring is a locally constant function on \(\mathbb Z_5\).
The augmentation splits by the measure at zero, so every continuous
functional on \(I\) extends continuously to that ring, and two extensions
differ by a constant. Each functional has finite five-power image
and factors through a finite group ring. The action on (D.75.8)
is the dual of the affine action (D.75.7). Its invariant classes
are therefore exactly the functions whose change under each
affine transformation is constant.

\paragraph{The exact global and local affine images.}
We prove this step for either \(k=\mathbb Q\) or \(k=\mathbb Q_5\).
The cyclotomic character maps \(G_k\) onto
\(\Gamma=\mathbb Z_5^\times\). For a nonzero class
\(b\in H^1(G_k,\mathbb Z_5(1))\) let \(d\) be its largest
divisibility exponent, and write \(b=5^d b_0\) in cohomology
with \(b_0\) primitive. After an integral coboundary change,
the actual cocycle is \(5^d\) times a cocycle for \(b_0\).

If \(\Delta_k=\ker\chi\), then
\[
                       b_0(\Delta_k)=\mathbb Z_5.          \tag{D.75.9}
\]
Indeed the Kummer groups here are torsion-free, so primitiveness
means that the image in \(H^1(G_k,\mathbb F_5(1))\) is nonzero.
The kernel of restriction to \(\Delta_k\) is
\(H^1(\Gamma,\mathbb F_5(1))=0\). For completeness, the
order-four subgroup of \(\Gamma\) acts by a nontrivial scalar
on \(\mathbb F_5(1)\); averaging and the procyclic resolution
of \(1+5\mathbb Z_5\) kill this first cohomology. Equivalently,
the cocycle equation for the commuting element with
cyclotomic value two makes every cocycle principal.
Thus \(b_0|_{\Delta_k}\) is nonzero modulo five. Its image
is a compact additive subgroup of \(\mathbb Z_5\), so is all
of \(\mathbb Z_5\). This proves (D.75.9).

It follows that the closed affine image, after the indicated
integral translation, is exactly
\[
                  5^d\mathbb Z_5\rtimes\mathbb Z_5^\times.
                                                               \tag{D.75.10}
\]
It contains all the displayed translations by (D.75.9).
For every multiplier its translation coordinate is in
\(5^d\mathbb Z_5\), and can be removed using one of those
translations. If the class is zero, its image after an integral
translation is instead exactly \(\mathbb Z_5^\times\), with
zero translation part. This is the \(d=\infty\) convention.
Thus (D.75.10) applies globally with \(d=r\) and locally with
\(d=N\). No finite-level approximation to the affine image
has been substituted here.

\paragraph{All invariant functions modulo constants.}
For finite \(d\), suppose \([f]\) in (D.75.8) is fixed by
the group in (D.75.10). For \(h\in5^d\mathbb Z_5\), write
\[
                         f(z+h)-f(z)=c(h).
 \]
The function \(c\) is an additive continuous character.
Pure multiplication by two belongs to (D.75.10), and
\(f(2z)-f(z)\) is constant; evaluating at zero shows that
this constant is zero. Consequently
\[
                         c(2h)=c(h)=2c(h),\qquad c(h)=0.
 \]
Thus \(f\) is genuinely periodic modulo \(5^d\), not just
periodic modulo constants. Similarly each pure unit multiplier
fixes \(f\) exactly, again by evaluating at zero.
Its values are precisely arbitrary values on the \(d+1\)
unit orbits of \(\mathbb Z/5^d\): the valuation shells
\(0,\ldots,d-1\) and the zero orbit. Modulo constants,
\[
 D_a^{G_k}\simeq(\mathbb Q_5/\mathbb Z_5)^d.
                                                               \tag{D.75.11}
\]
A concrete basis consists of the classes of the characteristic
functions of \(5^j\mathbb Z_5\), \(1\le j\le d\), with arbitrary
coefficients in \(\mathbb Q_5/\mathbb Z_5\).

When \(d=\infty\), invariance means that \(f\) is constant
on every valuation shell and at zero. Continuity into the
discrete target says that the shell values eventually equal
the value at zero. Hence
\[
 D_a^{G_k}\simeq
                   \bigoplus_{j\ge1}\mathbb Q_5/\mathbb Z_5.
                                                               \tag{D.75.12}
\]
The same characteristic functions form a basis, now with
finite support. These calculations include every invariant
class in (D.75.8), not only the selected ball functions.

We also identify the inclusion of global in local invariants.
Normalize the global cocycle to be divisible by \(5^r\).
A further integral translation \(d_0\) normalizes the local
cocycle to be divisible by \(5^N\), or to be zero when
\(N=\infty\). This translation satisfies \(d_0\in5^r\mathbb Z_5\):
evaluate the coboundary equation on an element with
\(\chi=2\). Both cocycle values being compared lie in
\(5^r\mathbb Z_5\), and \(\chi-1=1\).
Translation by \(d_0\) therefore leaves the first \(r\)
ball functions unchanged. Thus the inclusion identifies
the first \(r\) basis vectors in (D.75.11)--(D.75.12), and
\[
 \frac{D_a^{G_{\mathbb Q_5}}}{D_a^{G_{\mathbb Q}}}
 \simeq
 \begin{cases}
       (\mathbb Q_5/\mathbb Z_5)^{N-r},&N<\infty,\\
       \displaystyle\bigoplus_{\mathbb N}\mathbb Q_5/\mathbb Z_5,
                                                    &N=\infty .
 \end{cases}                                               \tag{D.75.13}
\]
In particular a global-invariant class cannot conceal a
nonzero ball function in this quotient. The preceding
translation argument explains the possible additive slope,
and the tame multiplier two forces that slope to vanish.

\paragraph{Passage to the actual compact global term.}
Put \(D_s=A_s^D\), a discrete union of the finite Cartier
duals used in D.63. The stable quotient \(A_s\twoheadrightarrow M_a\)
gives a genuine inclusion \(D_a\hookrightarrow D_s\).
Moreover
\[
 D_a^{G_{\mathbb Q_5}}\cap D_s^{G_{\mathbb Q}}
                         =D_a^{G_{\mathbb Q}}.
 \]
Consequently (D.75.13) injects into
\[
                      D_s^{G_{\mathbb Q_5}}/D_s^{G_{\mathbb Q}}.
                                                               \tag{D.75.14}
\]
Local duality and finite-level inverse limits identify the
Pontryagin dual of (D.75.14) with
\[
 \ker\{H^2_{\rm cts}(G_{\mathbb Q_5},A_s)
                      \longrightarrow\mathcal T_s\},
 \]
where \(\mathcal T_s\) is D.63's global-invariant dual.
Here invariants commute with the increasing union of finite
discrete duals, and all local finite cohomology groups are
finite. These are exactly the local duality and inverse-limit
identifications in D.63, based on Milne ADT I.2 and I.4
\cite{DE_MilneADT}. In particular no finite-level class
is being assumed to survive an inverse limit.

D.63 gives a continuous surjection from \(\mathcal K_s\)
onto this kernel. Dualizing the genuine inclusion (D.75.13)
and composing proves (D.75.2).
The splitting assertion also holds topologically. The discrete
dual of its target is a direct sum of copies of
\(\mathbb Q_5/\mathbb Z_5\), hence is divisible and injective
as a \(\mathbb Z_5\)-module. For instance the extension
criterion for the ideals \(5^n\mathbb Z_5\) is exactly
divisibility by \(5^n\), and the usual maximal-extension
argument extends a map from any submodule. Thus its inclusion
in \(\mathcal K_s^D\) has a retraction. Pontryagin duality
turns that retraction into a continuous section of (D.75.2).
No splitting of the nonseparated cohomology quotient is asserted.

\paragraph{Exact arithmetic parameters and examples.}
The completed Kummer group of \(\mathbb Q_5\) has the basis
\(\kappa(5),\kappa(6)\). Indeed
\[
 \mathbb Q_5^\times=5^{\mathbb Z}\times\mu_4\times(1+5\mathbb Z_5),
 \]
the prime-to-five factor disappears in the pro-five completion,
and logarithm identifies \(1+5\mathbb Z_5\) with \(5\mathbb Z_5\);
\(\log6\) has valuation one. For \(p\ne5\), its local
Kummer coordinate is exactly \(\lambda_p\) above, since
raising to the fourth power removes its Teichmuller part.
All \(\lambda_p\) are integral and \(a_p\to0\), so the
sum defining \(b_5(a)\) converges independently of order.
This proves (D.75.3).

For \(q\equiv1\pmod5\), \(v_5(\log q)=v_5(q-1)\);
this follows immediately from the logarithm series, whose
first term has strictly smaller valuation than every later
term. Thus \(v_5(q-1)=N+1\) gives exactly the stated local
divisibility \(N\). In particular \(101-1=4\cdot25\).

Finally \(\log16\) and \(\log81\) both have valuation one,
since \(16-1=15\) and \(81-1=80\). Their ratio is a
five-adic unit, and its residue is \(3/16=3\) in \(\mathbb F_5\).
The local unit coordinate of (D.75.4) is exactly zero,
while its coefficient at the prime two is one. Thus
\(N=\infty\) and \(r=0\), as claimed.
These are completed tangential parameters in the actual
zero-cusp packet, so their section fibres are nonempty
by D.35. The transformations permuting the punctures
transport the same proof and parameters to the other packets.

The executed certificate \texttt{code/de\_compact\_axis.py}
checks the affine unit-orbit decomposition and the independence
of all displayed ball functions at depths zero through five.
For the latter it evaluates their translation-one defects at
\(z=5^i-1\); the matrix has entries \(1\) for \(j\le i\)
and zero otherwise, hence determinant one. At \(z=1\)
all defects are zero, excluding an unnoticed constant defect.
It also computes (D.75.4)'s ratio modulo \(5^{20}\) with a
proved logarithm-tail bound and checks sample primes of exact
depths one through five. The output and JSON are retained.
The continuous duality and infinite-factor assertion follow
from the proof above, not from extrapolating these finite checks.
\end{proof}
